% !TEX root = ../main.tex
% =====================================================================
%  CHƯƠNG 6 — KẾT QUẢ MÔ PHỎNG TỔNG QUÁT
% ---------------------------------------------------------------------
%  Người phụ trách : [TODO: Họ tên - MSSV]
%  Tiền tố label   : c6      (fig:c6:..., tab:c6:..., eq:c6:..., sec:c6:...)
%  Thư mục ảnh     : image/chuong6/
%  File TLTK       : references/chuong6.bib   (key: c6_...)
%  Độ dài gợi ý    : 6–10 trang
%  Đáp ứng yêu cầu : (d) và kiểm chứng (e) trên toàn hệ thống.
%  Mục tiêu        : GHÉP 3 thành phần (Chương 3 + Chương 4 + Chương 5) thành
%                    một mạch duy nhất và mô phỏng ĐẦU CUỐI:
%                    khối lượng -> Load cell -> ... -> mã số ADC -> khối lượng tính lại.
%  Quy ước         : - KHÔNG lặp lại tính toán của Chương 3, 4, 5; chỉ tham chiếu
%                      bằng \ref{sec:c3:...}, \ref{tab:c5:...}.
%                    - Gom kết quả vào BẢNG TỔNG HỢP; KHÔNG tạo mỗi khối lượng
%                      một \section.
% =====================================================================
\chapter{Kết quả mô phỏng tổng quát}
\label{chap:c6}

% ---------------------------------------------------------------------
\section{Sơ đồ mô phỏng toàn hệ thống}
\label{sec:c6:so-do}

Mạch mô phỏng toàn hệ thống được ghép từ mạch đọc Load cell, chuyển đổi áp thành dòng và đọc dòng
(Chương~\ref{chap:c3}), mạch bù nhiệt (Chương~\ref{chap:c4}) và mạch đọc ADC (Chương~\ref{chap:c5}),
với tín hiệu giao tiếp giữa các khối theo Bảng~\ref{tab:c2:giao-tiep-khoi}.

\hinh[0.95\textwidth]{chuong6/so_do_toan_he_thong.png}{Sơ đồ mô phỏng toàn hệ thống}{fig:c6:so-do-toan-he-thong}

\TODO{Mô tả cách ghép nối, các điểm đo (test point) trên sơ đồ, những điều chỉnh khi ghép (nếu có).}

% ---------------------------------------------------------------------
\section{Kịch bản mô phỏng}
\label{sec:c6:kich-ban}

\TODO{Liệt kê các trường hợp mô phỏng: dải khối lượng (VD: 0, 2 mg, 100 g, 500 g, 1 kg) và dải nhiệt độ; phương pháp quy đổi mã ADC về khối lượng.}

Khối lượng tính lại từ mã số ADC được xác định theo:
\begin{equation}
    m_{\text{tính}} = \TODO{}
    \label{eq:c6:khoi-luong-tinh-lai}
\end{equation}

% ---------------------------------------------------------------------
\section{Kết quả mô phỏng theo khối lượng}
\label{sec:c6:theo-khoi-luong}

\TODO{Hình mô phỏng tiêu biểu (dùng \textbackslash haihinh), tại nhiệt độ chuẩn \qty{25}{\celsius}.}

\begin{table}[H]
    \centering
    \caption{Kết quả mô phỏng đầu cuối theo khối lượng tại \qty{25}{\celsius}}
    \label{tab:c6:theo-khoi-luong}
    \renewcommand{\arraystretch}{1.3}
    \resizebox{\textwidth}{!}{%
    \begin{tabular}{|c|c|c|c|c|c|c|}
        \hline
        \textbf{$m$ (\unit{\gram})} &
        \textbf{$V_{\text{KĐ}}$ (\unit{\volt})} &
        \textbf{$I_{\mathrm{lt}}$ (\unit{\milli\ampere})} &
        \textbf{$I_{\mathrm{mp}}$ (\unit{\milli\ampere})} &
        \textbf{Mã ADC} &
        \textbf{$m_{\text{tính}}$ (\unit{\gram})} &
        \textbf{Sai số (\unit{\percent})} \\ \hline
        0    & \TODO{} & 4.00  & \TODO{} & \TODO{} & \TODO{} & \TODO{} \\ \hline
        100  & \TODO{} & 5.60  & \TODO{} & \TODO{} & \TODO{} & \TODO{} \\ \hline
        500  & \TODO{} & 12.00 & \TODO{} & \TODO{} & \TODO{} & \TODO{} \\ \hline
        1000 & \TODO{} & 20.00 & \TODO{} & \TODO{} & \TODO{} & \TODO{} \\ \hline
    \end{tabular}}
\end{table}

\nhanxet{\TODO{Độ tuyến tính và sai số theo khối lượng.}}

% ---------------------------------------------------------------------
\section{Kết quả mô phỏng theo nhiệt độ}
\label{sec:c6:theo-nhiet-do}

\TODO{Mô phỏng toàn hệ thống ở các nhiệt độ khác nhau, có và không có mạch bù (so sánh với Bảng~\ref{tab:c4:ket-qua-bu-nhiet}).}

% ---------------------------------------------------------------------
\section{So sánh mô phỏng riêng lẻ và mô phỏng toàn hệ thống}
\label{sec:c6:so-sanh}

\TODO{So sánh với kết quả riêng lẻ ở Bảng~\ref{tab:c3:tong-ket} và Bảng~\ref{tab:c5:ket-qua-adc}; giải thích sai lệch phát sinh khi ghép (tải, offset tích lũy, ...).}

% ---------------------------------------------------------------------
\section{Đánh giá so với yêu cầu đề tài}
\label{sec:c6:danh-gia}

\begin{table}[H]
    \centering
    \caption{Đối chiếu kết quả mô phỏng với yêu cầu đề tài}
    \label{tab:c6:doi-chieu-yeu-cau}
    \renewcommand{\arraystretch}{1.3}
    \begin{tabularx}{\textwidth}{|l|c|c|X|}
        \hline
        \textbf{Yêu cầu} & \textbf{Mục tiêu} & \textbf{Đạt được} & \textbf{Nhận xét} \\ \hline
        Tầm đo          & \qtyrange{0}{1}{\kilo\gram}     & \TODO{} & \TODO{} \\ \hline
        Độ phân giải    & \qty{2}{\milli\gram}            & \TODO{} & \TODO{} \\ \hline
        Ngõ ra dòng     & \qtyrange{4}{20}{\milli\ampere} & \TODO{} & \TODO{} \\ \hline
        Sai số          & $\pm\qty{5}{\percent}$          & \TODO{} & \TODO{} \\ \hline
    \end{tabularx}
\end{table}

\nhanxet{\TODO{Kết luận hệ thống có đạt yêu cầu hay không, nguyên nhân sai số chính.}}
